% SOURCE_AUTHORITY: sga5_fr_workpass.tex, lines 8760--8795 (LF-normalized unit).
% SOURCE_SHA256: 791F4EFFC5E02832D5D77ED03518C8156D6F07E4C8238B03545DB93D883FBB28
\subsection*{2. Formalismo da cohomologia \(\ell\)-ádica}

\subsubsection*{2.1. Imagem inversa}

Seja
\[
f:X\longrightarrow Y
\]
um morfismo de pré-esquemas localmente noetherianos. O funtor imagem inversa
\[
f^*:\Ab(Y)\longrightarrow \Ab(X)
\]
é exato e, portanto, se prolonga de maneira evidente a um funtor exato, denotado da mesma forma,
\[
f^*:\Hom(\mathbb N^\circ,\Ab(Y))\longrightarrow \Hom(\mathbb N^\circ,\Ab(X)).
\]
Este último funtor transforma evidentemente um \(\mathbb Z_\ell\)-feixe construtível em um \(\mathbb Z_\ell\)-feixe construtível, e um sistema projetivo localmente \(AR\)-nulo em um sistema projetivo localmente \(AR\)-nulo. Com a notação de (1.5.4), induz, portanto, dois funtores imagem inversa
\[
f^*:\mathbb Z_\ell\text{-}\fc(Y)\longrightarrow \mathbb Z_\ell\text{-}\fc(X),
\]
\[
f^*:U(Y)\longrightarrow U(X),
\]
O segundo deles é exato.

A equivalência de categorias (1.5.5) e o fato de que \(f^*\) transforma um sistema projetivo localmente \(AR\)-nulo em um sistema projetivo localmente \(AR\)-nulo mostram então que o funtor imagem inversa
\[
f^*:\mathbb Z_\ell\text{-}\fc(Y)\longrightarrow \mathbb Z_\ell\text{-}\fc(X)
\]
É também exato.

Se \(g:Y\to Z\) é outro morfismo de pré-esquemas localmente noetherianos, os funtores \(f^*\) e \(g^*\) são evidentemente ligados por um isomorfismo
\[
(g\circ f)^*\xrightarrow{\sim}f^*\circ g^*,
\]
com a propriedade de associatividade habitual (SGA 1 VI 7 e ss.).
